\chapter{Problem Statement \& Desired Properties}
\label{ch:problem}

This chapter formalizes the network decision problems addressed in this thesis. We define the system model, state the operational decision spaces across single-path and multipath regimes, enumerate the essential properties a satisfactory solution must exhibit, and outline our core operating assumptions and trust boundaries.

\section{Problem Formulation}
\label{sec:problem:formulation}

Consider an endpoint host communicating over a path-aware architecture such as SCION. At any decision epoch $t$, the host discovers a dynamic set of $N_t$ candidate end-to-end paths $\mathcal{P}_t = \{p_1, \dots, p_{N_t}\}$. Each path $p_i$ is characterized by partially observable and dynamic performance properties -- latency, loss rate, available bandwidth, hop count, and path trust. Simultaneously, the running application expresses its \emph{intent} $g \in \mathcal{G}$, defined as a declarative preference vector prioritizing specific performance metrics (e.g., maximizing throughput, minimizing delay, or optimizing video quality).

Because network utilization and link availability fluctuate over time, the host must dynamically determine how best to utilize the path set $\mathcal{P}_t$ to serve intent $g$. We formulate this decision across two progressive operational settings:

\begin{enumerate}
  \item \textbf{Intent-Aware Single-Path Selection (Discrete Decision):} For general application flows constrained to a single path, the decision is a discrete selection $p^\star \in \mathcal{P}_t$. The objective is to maximize an intent-weighted reward balancing goodput, latency, and trust while minimizing the active probing overhead required to inspect candidate path states (\cref{ch:single_path}).
  
  \item \textbf{Joint Application--Network Multipath Optimization (Continuous Decision):} For high-bandwidth, latency-sensitive applications (e.g., real-time WebRTC video) operating over Multipath QUIC, the decision space expands. The endpoint must jointly choose a continuous traffic split $\boldsymbol{w} = (w_1, \dots, w_{N_t})$ across paths (where $\sum w_i = 1$) \emph{and} adapt the application's own sending rate (the video encoder's target bitrate $R_{\mathrm{video}}$) to maximize real-time Quality of Experience (QoE) (\cref{ch:multipath}).
\end{enumerate}

In both formulations, the technical crux is twofold: (i) the objective is dynamic and dictated by an intent $g$ that varies across flows, and (ii) the action space is dynamic and unbounded, as candidate path sets $\mathcal{P}_t$ fluctuate across network topologies and over time.

\section{Desired Properties}
\label{sec:problem:properties}

A satisfactory solution must exhibit six key properties, which also define the primary metrics evaluated in subsequent chapters (\cref{ch:single_path,ch:multipath,ch:synthesis}):

\begin{description}
  \item[Intent Alignment.] The policy must faithfully reflect the specified preference vector $g$. Decisions under $g_{\mathrm{latency}}$ must prioritize low delay, whereas decisions under $g_{\mathrm{throughput}}$ must prioritize goodput. A policy whose action choices remain invariant to changes in $g$ fails this requirement.
  
  \item[Zero-Shot Generalization Across Objectives.] A single trained network must serve the entire spectrum of application intents without requiring per-objective model retraining or fine-tuning.
  
  \item[Dynamic-Action Handling.] The policy architecture must natively accept dynamic candidate path counts $N_t$ without structural modification or retraining, remaining invariant to the arbitrary ordering of paths in $\mathcal{P}_t$.
  
  \item[Low Overhead.] Decision-making must remain computationally lightweight. In single-path selection, active probing cost must be minimized; in multipath real-time control, per-frame inference latency must fit comfortably within frame processing budgets ($\approx 33\,\mathrm{ms}$).
  
  \item[Structural \& Topological Generalization.] Performance must hold off-distribution across unseen topologies, dynamic path counts, and shifting traffic workloads.
  
  \item[Robustness to Heterogeneous Path Headroom.] The policy must sustain high performance when the
  load a naive allocation places on a path exceeds that path's residual capacity --- whether because
  the paths differ in capacity to begin with, or because path churn, link failures, and transient
  congestion bursts destroy an evenness that held a moment earlier. We state the property this way
  rather than as robustness to volatility because \cref{sec:p2eval:regimes} finds that headroom, not
  motion, is the governing variable: a stationary but asymmetric topology defeats a naive allocation as
  thoroughly as a churning one, and a stationary topology of near-equal paths leaves a scheduler almost
  nothing to earn.
\end{description}

For the single-path system, each of the six is measured rather than asserted:
\cref{sec:p1eval:intent} for Intent Alignment, \cref{sec:p1eval:zeroshot} for Zero-Shot Generalization,
\cref{sec:p1eval:dynamic} for Dynamic-Action Handling, \cref{sec:p1eval:overhead} for Low Overhead, and
\cref{sec:p1eval:ceiling} for behavior as the paths a flow can reach fill up, with
\cref{sec:p1eval:ablation} reporting Structural Generalization on a disjoint set of
source--destination pairs. One qualification applies throughout: structural generalization is
demonstrated across pairs and candidate-set sizes within a single topology, since only one topology was
generated, so generalization across topology \emph{scale} remains a design property rather than an
evaluated one.

The multipath system measures five of the same six properties, and the reader should have an index into
them: Intent Alignment (\cref{sec:p2eval:transfer,sec:p2eval:conditioning}), Zero-Shot Generalization
(\cref{sec:p2eval:conditioning}, with the caveat that the interior of the intent segment is
in-distribution there rather than unseen), Dynamic-Action Handling
(\cref{sec:p2eval:ablation,sec:p2eval:pathcount}), Structural Generalization (the held-out dynamics
arms of \cref{sec:p2eval:regimes,sec:p2eval:ablation}), and Robustness to Heterogeneous Path Headroom
(\cref{sec:p2eval:regimes,sec:p2eval:pathcount}). Low Overhead takes a different form in the two
systems: probing cost is the binding constraint for single-path selection
(\cref{sec:p1eval:overhead}), whereas for the multipath system it is per-frame inference latency,
reported in \cref{sec:p2eval:ablation,sec:p2eval:pathcount} and synthesized in
\cref{sec:analysis:overhead}.

\section{Assumptions \& Scope}
\label{sec:problem:assumptions}

We explicitly state the operating assumptions underlying our design:

\begin{enumerate}
  \item \textbf{Path-Aware Substrate.} The underlying network exposes multiple distinct end-to-end paths to the endpoint and honors source routing decisions, as provided by SCION. We do not assume control over internal intra-AS routing decisions.
  
  \item \textbf{Observable Path Feedback.} Per-path metrics (latency, loss, bandwidth) are measurable via active probing (in single-path selection) or derived continuously from transport feedback (e.g., QUIC ACK frames and RTT samples).
  
  \item \textbf{Explicit Intent Declaration.} Application intent is declared as an explicit preference vector (e.g., reward profile weights or target QoE parameterizations). We do not attempt to infer latent intent from raw packet inspection.
  
  \item \textbf{Episodic Stationarity.} Network state is assumed roughly stationary within a brief decision epoch, though conditions may drift or shift abruptly between epochs.
\end{enumerate}